\begin{tikzpicture}[x=0.043cm,y=3.0cm]
\draw[-{Latex[length=1.5mm]}] (-108,0) -- (108,0) node[below,font=\scriptsize] {$(V-U_j^{\,d})$ [mV]};
\draw[-{Latex[length=1.5mm]}] (-105,-0.02) -- (-105,1.02) node[above,font=\scriptsize] {$\xi_\eq$ \& scaled $\dd\xi_\eq/\dd V$};
\foreach \xx in {-100,-50,50,100} {\draw (\xx,0.012) -- (\xx,-0.012) node[below,font=\scriptsize] {\xx};}
\draw (0,0.012) -- (0,-0.012) node[below,font=\scriptsize] {0};
\draw (-108,0.5) -- (-102,0.5) node[left,font=\scriptsize] {0.5};
\draw (-108,1.0) -- (-102,1.0) node[left,font=\scriptsize] {1};
% ---- xi_eq(V) logistic, three T (thick = primary family) ----
\draw[thick] plot[smooth] coordinates {(-100,0.0130) (-90,0.0199) (-80,0.0303) (-70,0.0460) (-60,0.0693) (-50,0.1029) (-40,0.1503) (-30,0.2143) (-20,0.2961) (-10,0.3934) (0,0.5000) (10,0.6066) (20,0.7039) (30,0.7857) (40,0.8497) (50,0.8971) (60,0.9307) (70,0.9540) (80,0.9697) (90,0.9801) (100,0.9870)}; % T=268K, dash below applied via named style
\draw[thick,densely dashed] plot[smooth] coordinates {(-100,0.0130) (-90,0.0199) (-80,0.0303) (-70,0.0460) (-60,0.0693) (-50,0.1029) (-40,0.1503) (-30,0.2143) (-20,0.2961) (-10,0.3934) (0,0.5000) (10,0.6066) (20,0.7039) (30,0.7857) (40,0.8497) (50,0.8971) (60,0.9307) (70,0.9540) (80,0.9697) (90,0.9801) (100,0.9870)};
\draw[thick] plot[smooth] coordinates {(-100,0.0200) (-90,0.0292) (-80,0.0425) (-70,0.0615) (-60,0.0881) (-50,0.1249) (-40,0.1740) (-30,0.2372) (-20,0.3146) (-10,0.4039) (0,0.5000) (10,0.5961) (20,0.6854) (30,0.7628) (40,0.8260) (50,0.8751) (60,0.9119) (70,0.9385) (80,0.9575) (90,0.9708) (100,0.9800)};
\draw[thick,densely dotted] plot[smooth] coordinates {(-100,0.0282) (-90,0.0398) (-80,0.0557) (-70,0.0775) (-60,0.1069) (-50,0.1457) (-40,0.1954) (-30,0.2570) (-20,0.3301) (-10,0.4125) (0,0.5000) (10,0.5875) (20,0.6699) (30,0.7430) (40,0.8046) (50,0.8543) (60,0.8931) (70,0.9225) (80,0.9443) (90,0.9602) (100,0.9718)};
% ---- scaled derivative bell w|dxi/dV| (thin = secondary family), same T dash code ----
\draw[thin,densely dashed] plot[smooth] coordinates {(-100,0.0444) (-90,0.0675) (-80,0.1019) (-70,0.1521) (-60,0.2233) (-50,0.3199) (-40,0.4425) (-30,0.5833) (-20,0.7220) (-10,0.8267) (0,0.8661) (10,0.8267) (20,0.7220) (30,0.5833) (40,0.4425) (50,0.3199) (60,0.2233) (70,0.1521) (80,0.1019) (90,0.0675) (100,0.0444)};
\draw[thin] plot[smooth] coordinates {(-100,0.0609) (-90,0.0882) (-80,0.1267) (-70,0.1797) (-60,0.2504) (-50,0.3404) (-40,0.4477) (-30,0.5636) (-20,0.6717) (-10,0.7501) (0,0.7789) (10,0.7501) (20,0.6717) (30,0.5636) (40,0.4477) (50,0.3404) (60,0.2504) (70,0.1797) (80,0.1267) (90,0.0882) (100,0.0609)};
\draw[thin,densely dotted] plot[smooth] coordinates {(-100,0.0777) (-90,0.1081) (-80,0.1489) (-70,0.2024) (-60,0.2702) (-50,0.3522) (-40,0.4450) (-30,0.5405) (-20,0.6259) (-10,0.6859) (0,0.7076) (10,0.6859) (20,0.6259) (30,0.5405) (40,0.4450) (50,0.3522) (60,0.2702) (70,0.2024) (80,0.1489) (90,0.1081) (100,0.0777)};
% ---- central slope 1/(4w) tangent, T=298K reference (w=25.68 mV) ----
\draw[densely dotted,gray] (-20,0.3053) -- (20,0.6947);
\node[font=\scriptsize,align=left] at (46,0.60) {slope $\dd\xi_\eq/\dd V|_{0}$\\$=1/(4w)$};
% ---- compact legend (single block, minimal labels) ----
\node[font=\scriptsize,align=left,anchor=north west] at (-104,1.0) {\underline{$T$}: 298 K (solid) / 268 K (dashed) / 328 K (dotted)\\\underline{weight}: bold $=\xi_\eq$, thin $=$ scaled $\dd\xi_\eq/\dd V$};
\node[font=\scriptsize] at (0,-0.14) {center $U_j^{\,d}$};
\end{tikzpicture}
\caption{평형 진행률과 그 미분, 세 온도(식~\eqref{eq:xieq}$\cdot$eq.~\eqref{eq:belliden} — 좌표는 식 그대로의
수치 평가, 검산 = T6\_calc.py, $n_j{=}1$). 굵은 선 $=\xi_\eq$ logistic, 가는 선 $=$ 규격화한 미분
$w_j|d\xi_\eq/dV|=\xi_\eq(1-\xi_\eq)$(공통 배율로 재척도, 온도 간 상대 높이는 보존) — 선 굵기가 두 곡선군을,
선 종류(실선/파선/점선)가 $T{=}298/268/328$ K 를 가른다. 실선축 $x$ 는 이제 정규화된 $z$ 가 아니라 실제
$(V-U_j^{\,d})$[mV] 라, $w_j=n_jRT/F$ 가 온도에 따라 커지며(저온일수록 좁고 높은 종, 고온일수록 넓고 낮은
종) 곡선이 실제로 벌어지는 것이 보인다. 회색 점선은 $z{=}0$(중심)에서의 접선 — 기울기가 정확히
$1/(4w_j)$ 인 종의 기하적 의미(폭의 역수 $1/4$ 배가 곧 중심 기울기)다. 이 종이 다음 절의 평형 peak 이고,
방향에 불변(양수)이다.}
\label{fig:logistic}
